\documentclass[11pt, a4paper]{article}
\usepackage[utf8]{inputenc}
\usepackage{amsmath, amssymb}
\usepackage{graphicx}
\usepackage[margin=1in]{geometry}
\usepackage{tcolorbox}
\usepackage{fancyhdr}
\usepackage{hyperref}
\usepackage{listings}
\usepackage{xcolor}

\definecolor{UNIBBlue}{RGB}{0, 51, 102}
\definecolor{codegreen}{rgb}{0,0.6,0}
\definecolor{codegray}{rgb}{0.5,0.5,0.5}
\definecolor{codepurple}{rgb}{0.58,0,0.82}
\definecolor{backcolour}{rgb}{0.95,0.95,0.92}

\lstdefinestyle{mystyle}{
    backgroundcolor=\color{backcolour},   
    commentstyle=\color{codegreen},
    keywordstyle=\color{magenta},
    numberstyle=\tiny\color{codegray},
    stringstyle=\color{codepurple},
    basicstyle=\ttfamily\footnotesize,
    breakatwhitespace=false,         
    breaklines=true,                 
    captionpos=b,                    
    keepspaces=true,                 
    numbers=left,                    
    numbersep=5pt,                  
    showspaces=false,                
    showstringspaces=false,
    showtabs=false,                  
    tabsize=2
}
\lstset{style=mystyle}

\pagestyle{fancy}
\fancyhf{}
\rhead{\textbf{Optimisasi untuk Teknik Elektro}}
\lhead{Lembar Kerja Mahasiswa (LKM) - Minggu 5}
\cfoot{\thepage}

\begin{document}

\begin{center}
    \Large\color{UNIBBlue}\textbf{LEMBAR KERJA MAHASISWA (LKM)}\\
    \large\textbf{Minggu 5: Implementasi Algoritma Berbasis Gradien}
\end{center}

\vspace{0.5cm}
\textbf{Nama Praktikan :} \rule{6cm}{0.4pt} \\
\textbf{NPM :} \rule{6.8cm}{0.4pt}

\section*{Tujuan Praktikum}
1. Mampu mengimplementasikan algoritma Gradient Descent dari awal (from scratch). \\
2. Mampu menganalisis pengaruh \textit{learning rate} terhadap konvergensi.

\section*{Kasus Uji: Meminimalkan Fungsi 2D Sederhana}
Misalkan kita memiliki fungsi biaya operasional sistem:
\[ f(x, y) = x^2 + 2y^2 + 2x - 4y + 5 \]

\section*{Aktivitas 1: Turunan Parsial (C3)}
Tentukan turunan parsial terhadap $x$ dan $y$:
\begin{align*}
    \frac{\partial f}{\partial x} &= \ldots \\
    \frac{\partial f}{\partial y} &= \ldots
\end{align*}

\section*{Aktivitas 2: Implementasi Python (C4)}
Ketik dan jalankan kode berikut di Jupyter Notebook/Google Colab:

\begin{lstlisting}[language=Python]
import numpy as np

# Definisikan fungsi dan gradien
def f(x, y):
    return x**2 + 2*y**2 + 2*x - 4*y + 5

def grad_f(x, y):
    df_dx = 2*x + 2
    df_dy = 4*y - 4
    return np.array([df_dx, df_dy])

# Hyperparameters
alpha = 0.1 # Learning rate
epochs = 50
p = np.array([5.0, 5.0]) # Tebakan awal (Initial guess)

history = []

for i in range(epochs):
    gradient = grad_f(p[0], p[1])
    p = p - alpha * gradient
    history.append(f(p[0], p[1]))
    if i % 10 == 0:
        print(f"Iterasi {i}: Koordinat = {p}, f(x,y) = {f(p[0], p[1]):.4f}")

print(f"Optimal di temukan di: {p}")
\end{lstlisting}

\section*{Aktivitas 3: Analisis Konvergensi (C5)}
Ubah nilai `alpha` pada program di atas menjadi:
a) `alpha = 0.01` \\
b) `alpha = 0.5` \\
c) `alpha = 1.0` \\

Amati keluaran setiap skenario. Apa yang terjadi pada nilai koordinat dan $f(x,y)$ pada setiap kondisi?

\vspace{1cm}
\textbf{Laporan Pengamatan:} \\
\rule{\textwidth}{0.4pt} \\
\rule{\textwidth}{0.4pt} \\
\rule{\textwidth}{0.4pt} \\
\rule{\textwidth}{0.4pt}

\end{document}
